% !TEX TS-program = pdflatex
% !TeX program = pdflatex
% !TEX encoding = UTF-8
% !TeX spellcheck = en_US

\documentclass[11pt, a4paper]{article}
%\usepackage{fullpage}
\usepackage[left=1cm,right=1cm,top=1cm,bottom=2cm]{geometry}
\usepackage[fleqn]{amsmath}
\usepackage{amssymb}
%\usepackage{indentfirst}
\usepackage[T1]{fontenc}
\usepackage[utf8]{inputenc}
\usepackage[french,english]{babel}
\usepackage{txfonts} 
\usepackage[]{graphicx}
\usepackage{multirow}
\usepackage{hyperref}
\usepackage{parskip}
\usepackage{multicol}
\usepackage{wrapfig}
\usepackage{amssymb}
\usepackage{longtable}
\usepackage{wasysym}
\usepackage{soul}


\usepackage{turnstile}%Induction symbole

\usepackage{tikz}
\usetikzlibrary{arrows, automata}
\usetikzlibrary{decorations.pathmorphing}

\renewcommand{\baselinestretch}{1}

\setlength{\parindent}{24pt}


\begin{document}

%\selectlanguage {french}
%\pagestyle{empty} 

\noindent
\begin{tabular}{ll}
\multirow{3}{*}{\includegraphics[width=1.5cm]{../extra/logo/esi.nlp.pdf}} & 
\'Ecole national Supérieure d'Informatique\\
& 2\textsuperscript{nd} year second cycle (2CSSID)\\
& NLP: Natural Language Processing (2022-2023)
\end{tabular}\\[.25cm]
\noindent\rule{\textwidth}{2pt}\\[-0.5cm]
\begin{center}
{\LARGE \textbf{Tutorial 03: Sentence semantics and discourse}}
\begin{flushright}
	Abdelkrime Aries
\end{flushright}
\end{center}\vspace{-0.5cm}
\noindent\rule{\textwidth}{2pt}

\section*{1. Sentence semantics}

Given these sentences:
\begin{center}
	\begin{tabular}{|l|}
	\hline 
	\\
	\ul{Each student} prepares \ul{a report} with \ul{\LaTeX} \\
	If \ul{a student} is subscribed to \ul{ESI}, so \ul{he} studies \ul{CS}\\
	\\
	\hline 
\end{tabular}
\end{center}

\begin{enumerate}
	\item For each underlined expression, indicate the thematic role using the table seen in class.
	\item Suppose we used an LSTM to estimate these roles, write the result of each sentence using IOB representation.
	\item Using FrameNet (\url{https://framenet.icsi.berkeley.edu/fndrupal/} [visited 2023-05-23]), indicate the activated semantic frameworks for each sentence.
	\item Using PropBank (\url{https://propbank.github.io/} [visited 2023-05-23]), indicate the activated predicates.
	\item Annotate the two sentences based on the arguments of each predicate.
	\item Provide a context-free grammar to generate these sentences.
	\item Represent the two sentences in first-order logic where the domain provides the following predicates:
	Report(r); Student(s); Prepare(x, y); Prepare-With(x); Subscribe-To(x, y); Studies(x, y); CS(c).
	It also provides the following constants: ESI; LATEX.
	\item Enrich the previous grammar with semantic rules in order to generate these representations.
	\item Draw the semantic tree of each sentence.
	\item Now we want to use a neo-Davidsonian representation.
	The domain, in this case, provides the following predicates:
	Student(s); Prepare(e); Report(r); Prepare-Agent(e, x); Prepare-Theme(e, x); Prepare-Instrument(e, x); Location(e, x); Subscribe(e); Study(e); CS(x).
	It also provides the following constants: ESI; LATEX.
	\item Enrich the previous grammar with semantic rules in order to generate these representations.
	\item Draw the semantic tree of each sentence.
\end{enumerate}

\section*{2. Coreference resolution}

Given this text:
\begin{center}
	\begin{tabular}{|p{0.8\textwidth}|}
	\hline
	\\
	A girl picked a flower.
	She felt it.
	It has a very good smell.
	This beautiful flower is in école nationale supérieure d'informatique where the girl is studying.
	ESI is located in Algiers; a city in northern Algeria.\\
	\\
	\hline
\end{tabular}
\end{center}

\begin{enumerate}
	\item Find all mentions indicating the form of reference.
	\item Annotate co-references.
	\item Write a rule-based function that takes two mentions as arguments and returns a boolean indicating whether the first is an antecedent of the second based on the properties of both.
	\item Apply Mention-Pair algorithm based on this function to the mentions of the first question.
	Then, check if the generated co-references are similar to those manually annotated.
	\item Write another function that takes two mentions with an array of mentions between them as arguments and returns a score indicating whether the first is an antecedent of the second based on the properties of both.
	\item Apply Mention-Rank algorithm based on this function to the mentions of the first question.
	Then, check if the generated co-references are similar to those manually annotated.
	\item Write a recursive function that implements an Entity-based model. 
	This function takes one set of clusters as an argument and returns another set of clusters.
\end{enumerate}

\section*{3. Coherence}

Voici un texte : 

\begin{center}
	\begin{tabular}{|p{0.8\textwidth}|}
	\hline
	\\
	Winter is one of the four seasons of the year. The report was well presented.
	Rain is one of the characteristics of this season.
	I read a report on this season. \\
	\\
	\hline 
\end{tabular}
\end{center}

\begin{enumerate}
	\item Represent each sentence using TF (eliminate stop words).
	\item Calculate cosine similarities between all sentences.
	\item Attempt to rearrange sentences by maximizing the sum of similarities between consecutive sentences.
	\item Did you find the best order?
	\item Based on the correct order, find the RST relations by drawing RST parse tree.
	\item Apply ``Centering theory" on mis-ordered and correct text.
	\item Suppose the CONTINUE, RETAIN, Smooth-SHIFT and Rough-SHIFT transitions are scored as 3, 2, 1 and 0 respectively.
	Calculate the score for each text based on the result of the previous question.
	\item Using ``Entity Grid model", represent the sentences by entities and by the probabilities of the grammatical transitions for the two texts: wrong and correct.
\end{enumerate}

\end{document}
